Designing an effective bike-sharing station network is one of the most important strategic decisions in the planning of docked \acp{bss}. Station placement affects system accessibility, influences potential demand, and directly affects management operations such as rebalancing. As reviewed in Chapter~\ref{ch:station_placement}, research in this field has progressed from early heuristic and \ac{gis}-based suitability analyses that relied on low-resolution spatial data to more advanced multi-criteria and optimization-driven approaches that make use of high-resolution datasets. Despite these advances, many existing methods remain limited in flexibility and are seldom adapted to the fine-grained urban realities of specific cities.

Part~\ref{pt:data-driven-station-planning} addresses this gap by developing and applying a flexible data-driven framework for station planning in Barcelona. After introducing the city in Chapter~\ref{ch:use_case_barcelona}, this part presents in Chapter~\ref{ch:station_placement} the modelling approach, which integrates high-resolution spatial data, topography-adjusted cycling distances, and network-based system accessibility and inter-station proximity metrics into a multi-objective optimization framework solved through a \ac{ga}. That chapter also illustrates how the framework supports scenario exploration and incremental network expansion using Barcelona as the case study. These chapters are adapted from Grau-Escolano et al.~\cite{GrauEscolano2025}, currently under review.